\chapter{Matières premières et énergie}\label{ch:iv2:commodities-and-energy}

\begin{ivtrains}
Les formes de courbe et ce qu'elles disent du stockage, les écarts entre produits et entre lieux, les courbes saisonnières et l'électricité, et les options qu'un actif physique représente réellement. Chaque niveau ici est inventé pour le calcul, sauf lorsqu'un encadré donne une date.
\end{ivtrains}

\begin{question}[$\star$ \iqroles{Trader, QT, Analyste} \iqfirm{toutes} \iqlevel{junior}]\label{iq:iv2:commodities-and-energy:1}
Le pétrole brut vaut 70 sur le marché spot, et le contrat à terme à un an vaut 74. Comment appelle-t-on cette forme de courbe~? Vous achetez le contrat à terme. Si le cours spot vaut toujours 70 dans un an, que gagnez-vous ou perdez-vous par baril~?
\end{question}

\begin{question}[$\star$ \iqroles{Trader, QT, QR} \iqfirm{toutes} \iqlevel{junior}]\label{iq:iv2:commodities-and-energy:2}
Une matière première vaut 80 aujourd'hui, le taux est de 5~\%, et le stockage coûte 4 pour l'année, payés en fin de période. Quel est le niveau maximal du contrat à terme à un an~? Il se négocie à 92, et vous avez accès au stockage. Que faites-vous, et combien verrouillez-vous~?
\end{question}

\begin{question}[$\star$ \iqroles{Trader, Analyste, QT} \iqfirm{hedge fund, banque} \iqlevel{junior}]\label{iq:iv2:commodities-and-energy:3}
Le brut vaut 80~\$ le baril, l'essence 2.40~\$ le gallon, le fioul domestique 2.70~\$ le gallon, et un baril contient 42 gallons. Quel est le crack spread 3:2:1 par baril de brut~? Qui est naturellement long de ce crack spread~?
\end{question}

\begin{question}[$\star$ \iqroles{Trader, QT} \iqfirm{hedge fund, banque} \iqlevel{junior}]\label{iq:iv2:commodities-and-energy:4}
Le brut de l'échéance la plus proche vaut 72, et le mois suivant vaut 70. Chaque contrat porte sur 1\,000 barils. Vous achetez l'échéance la plus proche et vendez le mois suivant. Que gagnez-vous si les deux montent de 5~\$~? Si la première monte de 3~\$ et la seconde de 2~\$~? Quel point de vue exprimez-vous~?
\end{question}

\begin{question}[$\star$ \iqroles{Trader, Analyste, QT} \iqfirm{banque, hedge fund} \iqlevel{junior}]\label{iq:iv2:commodities-and-energy:5}
L'électricité s'échange à 60~\$ le MWh et le gaz à 5~\$ le MMBtu, et une centrale a un taux de chaleur de 7 MMBtu par MWh. Quel est le spark spread~? Combien gagne une centrale de 100 MW sur un bloc de 16 heures~? Avec un prix du carbone de 30~\$ la tonne et 0.4 tonne par MWh, quel est le clean spark spread~?
\end{question}

\begin{question}[$\star\star$ \iqroles{QR, Trader, QT} \iqfirm{toutes} \iqlevel{junior, confirmé}]\label{iq:iv2:commodities-and-energy:6}
Un métal vaut 70 au comptant et le contrat à terme à un an vaut 66. Le taux est de 4~\% et le stockage coûte 2~\% de la valeur par an, les deux en capitalisation continue. Quel rendement d'opportunité cela implique-t-il~? Que signifie ce nombre, et quand serait-il élevé~?
\end{question}

\begin{question}[$\star\star$ \iqroles{Trader, QR, QT} \iqfirm{toutes} \iqlevel{junior, confirmé}]\label{iq:iv2:commodities-and-energy:7}
Le spot vaut 50, le taux est de 4~\% et le stockage coûte 1~\$ par mois. Quelle peut être l'ampleur du contango sur un mois~? Existe-t-il une limite équivalente à l'ampleur de la backwardation~? Pourquoi, ou pourquoi pas~?
\end{question}

\begin{question}[$\star\star$ \iqroles{QR, QT, Trader} \iqfirm{hedge fund, banque} \iqlevel{junior, confirmé}]\label{iq:iv2:commodities-and-energy:8}
La courbe du actif A monte de 1~\% par mois, indéfiniment, et son cours spot ne bouge jamais. Le gaz a une courbe purement saisonnière~: chaque mois calendaire a un prix de $3 + 0.6\cos(2\pi(m - 1)/12)$, et chaque contrat expire exactement à ce prix. Que rapporte sur un an une position longue reroulée chaque mois dans chacun des deux cas~? Quelle part de la pente de la courbe du gaz est un rendement réel~?
\end{question}

\begin{question}[$\star\star$ \iqroles{Trader, QT, QR} \iqfirm{hedge fund, banque} \iqlevel{junior}]\label{iq:iv2:commodities-and-energy:9}
Une compagnie aérienne doit acheter dans trois mois 10 millions de gallons de carburant d'aviation et se couvre avec des contrats à terme sur fioul domestique de 42\,000 gallons chacun. La volatilité mensuelle des variations du prix du carburant d'aviation est de 3~\%, celle du fioul domestique de 3.5~\%, et leur corrélation vaut 0.8. Quel ratio de couverture et combien de contrats~? Quelle part de la variance ce dispositif élimine-t-il~?
\end{question}

\begin{question}[$\star\star$ \iqroles{QR, Trader} \iqfirm{hedge fund, banque} \iqlevel{confirmé}]\label{iq:iv2:commodities-and-energy:10}
Un contrat à terme dont le temps jusqu'à l'échéance est $\tau$ a une volatilité de $0.6\,e^{-2\tau}$. Quelle est aujourd'hui la volatilité du contrat à terme à un mois, et celle du contrat à terme à un an~? Quelle volatilité moyenne utiliseriez-vous pour une option à un an sur le contrat à terme à un an~? Pour une option à un mois sur le contrat à terme à un mois~? Quelle est l'histoire économique~?
\end{question}

\begin{question}[$\star\star$ \iqroles{QT, QR, Trader} \iqfirm{hedge fund, banque} \iqlevel{junior, confirmé}]\label{iq:iv2:commodities-and-energy:11}
Une option à un mois à la monnaie sur un contrat à terme dont le cours vaut 70, avec une volatilité de 40~\%, se règle contre la moyenne du cours sur tout le mois. Quelle est sa valeur approximative, et que vaut une option européenne classique~? Pourquoi tant d'options sur matières premières se règlent-elles de cette façon~?
\end{question}

\begin{question}[$\star\star$ \iqroles{Trader, Analyste, QT} \iqfirm{hedge fund, banque} \iqlevel{junior, confirmé}]\label{iq:iv2:commodities-and-energy:12}
Le brut américain coûte 73 au port. Une cargaison met 30 jours pour atteindre l'Europe, et le fret coûte 3.50~\$ le baril. Le prix de référence européen pour une livraison à 30 jours vaut 77. Le financement coûte 5~\%. L'arbitrage est-il ouvert~? Comment le verrouillez-vous, et quel risque demeure-t-il~?
\end{question}

\begin{question}[$\star\star$ \iqroles{Trader, QT, Analyste} \iqfirm{banque, hedge fund} \iqlevel{junior, confirmé}]\label{iq:iv2:commodities-and-energy:13}
Un marché de l'électricité compte quatre ensembles de centrales~: éolien 30 MW à 0~\$, nucléaire 40 MW à 10~\$, cycle combiné au gaz 40 MW à 40~\$ et pics 20 MW à 90~\$. Quel est le prix pour une demande de 20, 60, 100 et 115 MW~? Que se passe-t-il à 140~? La capacité éolienne double pour atteindre 60 MW. Qu'advient-il des prix~? Pourquoi les prix de l'électricité peuvent-ils devenir négatifs~?
\end{question}

\begin{question}[$\star\star\star$ \iqroles{QR, Trader} \iqfirm{banque, hedge fund} \iqlevel{confirmé}]\label{iq:iv2:commodities-and-energy:14}
Le contrat à terme sur l'électricité vaut 40 par MWh et celui sur le gaz 5 par MMBtu. Une centrale a un taux de chaleur de 7 MMBtu par MWh. La volatilité de l'électricité est de 60~\%, celle du gaz de 40~\%, leur corrélation vaut 0.5, et l'option de fonctionnement expire dans trois mois. En négligeant l'actualisation, évaluez avec Margrabe l'option de fonctionnement pour un MWh. Quelle est la valeur intrinsèque, et que représente le reste~? Que se passe-t-il pour une corrélation de 0.9~? Comment la couvrez-vous~?
\end{question}

\begin{question}[$\star\star\star$ \iqroles{Trader, QR, Risque} \iqfirm{hedge fund, banque} \iqlevel{junior, confirmé}]\label{iq:iv2:commodities-and-energy:15}
Un contrat à terme sur pétrole brut à livraison physique n'a plus qu'un jour à courir. Le mois suivant se négocie à 20. Le stockage au point de livraison est pratiquement plein. Qu'est-ce qui empêche le contrat arrivant à échéance de se négocier sous zéro~? À quel coût de stockage, par baril et par mois, un prix de $-38$ a-t-il un sens~? Qui doit vendre à n'importe quel prix~?
\end{question}

\begin{question}[$\star\star\star$ \iqroles{QR, Trader, QD} \iqfirm{banque, hedge fund} \iqlevel{junior, confirmé}]\label{iq:iv2:commodities-and-energy:16}
Un réservoir de gaz d'une unité peut injecter ou soutirer une unité par mois. La courbe à terme sur cinq mois vaut 3.00, 2.90, 2.85, 3.40, 3.60, et le réservoir doit être vide à la fin. L'injection coûte 0.02 et le soutirage 0.01. Quelle est la valeur intrinsèque~? Puis le prix du mois 5 est révélé au mois 4, à 4.10 ou 3.10 avec des chances égales. Que vaut le réservoir si vous pouvez soutirer au mois 4 plutôt que d'attendre~?
\end{question}

\begin{question}[$\star\star\star$ \iqroles{Trader, Risque, QR} \iqfirm{banque, hedge fund} \iqlevel{confirmé}]\label{iq:iv2:commodities-and-energy:17}
Vous avez vendu à prix fixe 10 millions de barils, soit 1 million de barils par an pendant dix ans, et vous les couvrez un pour un avec des contrats à terme longs sur l'échéance la plus proche, reroulés chaque mois. Le pétrole baisse de 20~\$ dans un mouvement parallèle sur toute la courbe, et le taux est de 5~\%. Quels sont votre appel de marge et la valeur actuelle de votre gain sur les ventes~? Cette couverture est-elle juste~? Supposons maintenant que la courbe ne bouge que de 30~\% sur la partie longue, ou qu'elle passe en contango de 50 cents par mois. Qu'en est-il alors~?
\end{question}

\begin{question}[$\star\star\star$ \iqroles{Trader, QR, Analyste} \iqfirm{banque, hedge fund} \iqlevel{junior, confirmé}]\label{iq:iv2:commodities-and-energy:18}
Un producteur veut se protéger sur un contrat à terme à un an à 70, sans prime, en achetant un put à 60 et un call. Les taux sont nuls et la volatilité est plate à 30~\%. Quel strike de call rend l'opération sans coût~? Avec un skew, le put est à 36~\% de volatilité et le call à 26~\%. Qu'en est-il alors~? Que abandonne le producteur~?
\end{question}

\begin{question}[$\star\star\star$ \iqroles{QR, Trader} \iqfirm{banque, hedge fund} \iqlevel{confirmé}]\label{iq:iv2:commodities-and-energy:19}
Vous évaluez une option à six mois sur un spread de raffinage de gain $\max(F_1 - F_2 - 15, 0)$ avec $F_1 = 100$, $F_2 = 80$, des volatilités de 35~\% et 30~\%, et une corrélation de 0.9. Utilisez l'approximation de Kirk. Donnez ensuite une seconde estimation en traitant le spread comme gaussien. Laquelle retiendriez-vous, et comment la vérifieriez-vous~?
\end{question}

\begin{question}[$\star\star\star\star$ \iqroles{QR, Trader} \iqfirm{banque, hedge fund} \iqlevel{confirmé}]\label{iq:iv2:commodities-and-energy:20}
On vous propose un contrat de tolling pour un mois~: vous payez une redevance et pouvez faire fonctionner une centrale de 100 MW, taux de chaleur 8, pendant 22 jours à raison de 16 heures par jour, en achetant du gaz et en vendant de l'électricité. Le contrat à terme sur l'électricité vaut 42~\$, celui sur le gaz 5~\$. Les volatilités mensuelles à terme sont de 50~\% et 35~\%, la corrélation de 0.6, et le mois est à un mois. Les volatilités à un jour sont de 90~\% et 50~\%, la corrélation de 0.5, et chaque option quotidienne est traitée comme expirant dans $1/8$ d'année. Quelle est la valeur intrinsèque~? Que vaut-elle comme une seule décision mensuelle, et comme 22 décisions quotidiennes~? Quelle couverture mettez-vous en place contre chacune~? Quelle redevance paieriez-vous, et pourquoi pas le haut de votre fourchette~? Que devient la valeur si le contrat à terme vaut 32~\$~?
\end{question}
